\documentclass[11pt]{article}
\usepackage{icmlko}
\begin{document}

\title{감사가 가리킨 자리가 고칠 자리는 아니었다 --- 기각 하나가 감사에 문턱을 달아 줬다}
\author{월드모델 실험실}
\date{노트 593 · 2026년 8월 5일}
\maketitle

\begin{abstract}
매 사이클 감사(\texttt{lab.sideaudit})의 검사 ③ 은 ``학습과 유보의 결측률이
다른가, 그 마스크가 라벨과 붙나''를 \textbf{곱}으로 줄 세운다. 미처리
$1$ 위가 \textbf{웹툰 검색}이었다 --- 학습 덮음 $70\%$ 대 유보 $99\%$,
마스크 $\leftrightarrow$ 라벨 $+0.1471$, 곱 $0.043$.

팔 둘로 돌렸다. \textbf{둘 다 기각이다.}

\begin{center}
\begin{tabular}{lrrl}
\hline
자 & B(축 통째) & C(관측-$0$ 만) & \\
\hline
판 & $\mathbf{-0.0105}$ & $\mathbf{-0.0087}$ & 둘 다 $0/12$ \\
날짜 통제 판 & $-0.0090$ & $-0.0082$ & 둘 다 $0/12$ \\
무리 안 판 & $-0.0086$ & $-0.0115$ & \\
KR 만화 & $+0.0086$ & $+0.0045$ & \\
비게임 앱 & $-0.0029$ & $+0.0005$ & B 는 \textbf{거울 갈림} \\
\textbf{웹툰} & $\mathbf{-0.0284}$ & $\mathbf{-0.0260}$ & 둘 다 $0/12$ \\
\hline
\end{tabular}
\end{center}

\textbf{마스크가 수집 순서를 나르는 것은 맞다. 그런데 축이 주는 것이 더
크다.}
\end{abstract}

\begin{figure}[t]
\centering
\includegraphics[width=\textwidth]{figs/f1.pdf}
\caption{\textbf{가} --- 세 번 돌려서 문턱을 얻었다.
\textbf{나} --- 판이 내려가고 웹툰이 잃는다.
\textbf{다} --- 왜 웹툰만 다른가.}
\end{figure}

\section{왜 웹툰만 다른가}

같은 검사가 낸 후보를 이 판은 이제 \textbf{세 번} 돌려 봤다.

\begin{center}
\begin{tabular}{lrrrl}
\hline
도메인 & 곱 & 학습 $\to$ 유보 덮음 & 차 & 결과 \\
\hline
애니 & $0.200$ & $56\% \to 97\%$ & $41$\%p & 채택(노트 $498 \cdot 513$) \\
만화 & $0.072$ & $17\% \to 3\%$ & $14$\%p & 채택(노트 $553$) \\
\textbf{웹툰} & $\mathbf{0.043}$ & $\mathbf{70\% \to 99\%}$ & $29$\%p & \textbf{기각} \\
\hline
\end{tabular}
\end{center}

\begin{quote}
만화는 \textbf{유보에서 $3\%$} 라 축이 거의 쓸모없었다 --- 떼는 것이
이득이다. 애니는 학습 $56\%$ 에서 유보 $97\%$ 로 \textbf{차가 $41$\%p} 다.
\textbf{웹툰은 양쪽 다 높다}($70/99$) --- \textbf{학습에서도 대부분
관측되므로 축이 실제로 정보를 준다.}
\end{quote}

\textbf{차가 크다고 뗄 것이 아니다.} 차가 크면서 \emph{한쪽이 낮아야}
뗄 값이 있다.

\section{기각이 문턱을 줬다}

남은 미처리는 모바일 $0.031$ $\cdot$ 세계애니 $0.010$ $\cdot$ 도서
$0.003$ --- \textbf{전부 웹툰($0.043$)보다 아래다.} 웹툰이 기각이면
그 아래는 안 재도 된다.

\begin{center}
\texttt{sideaudit.CHASE\_HIT = 0.05}
\end{center}

감사가 이제 남은 것이 전부 문턱 아래면 \textbf{``안 쫓는다''고 직접
말한다} --- 그리고 그 근거(웹툰을 실제로 돌려 기각한 수치)를 같이 찍는다.

\begin{quote}
\textbf{표본이 셋뿐이라 문턱 자체는 약하다.} $0.043$ 과 $0.072$ 사이
어디인지는 모른다. 새 후보가 그 사이에 오면 그때 다시 본다 ---
\textbf{주석에 그렇게 적었다.}
\end{quote}

\section*{감사 ③ 이 비었다}

이 판에서 ``우리 쪽 사정이 축에 실렸나''(노트 $546 \cdot 552 \cdot 554$)는
가장 값비싼 갈래였다 --- 이번 사이클의 정정 $0.0246$ 이 검색 축 전체의
값($0.0204$)보다 컸다. 그 갈래의 세 번째 검사가 \textbf{이제 비었다.}

\begin{itemize}
\item ① 사후 상태를 읽는 축 --- \textbf{없음}
\item ② 시작일이 물리적으로 불가능한 행 --- 웹툰 $64$ 행(이미
\texttt{datehygiene} 로 처리) \emph{하나뿐}
\item ③ 수집 순서가 마스크에 실린 축 --- \textbf{곱 $0.05$ 위가 없다}
\item ④ 공유 축의 부호 --- \textbf{원 방향에서 갈리는 축 $0$ 개}
\end{itemize}

\begin{quote}
넷이 다 비었다는 것은 \emph{깨끗하다}가 아니라 \textbf{이 자들로는 더
못 찾는다}는 뜻이다. 다음 누출 층은 \textbf{새 자를 만들어야} 보인다 ---
이 판이 지금까지 네 층을 그렇게 찾았다(노트 $117$ 둘 $\cdot$ $141$ $\cdot$
$552$).
\end{quote}

\section*{조항}

\begin{quote}
\textbf{감사가 낸 순위는 후보 목록이지 할 일 목록이 아니다.} 곱이 크다고
고칠 것이 아니고, 실제로 돌려 봐야 안다. \textbf{그리고 돌려 본 결과를
감사에 되먹여야} 다음 사람이 같은 것을 다시 돌리지 않는다 --- 노트 $561$ 이
``이미 고친 것''을 표시하게 한 것과 같은 자리이고, 여기서는
``\textbf{돌려 봤는데 아니었던 것}''이다.
\end{quote}

\section*{남은 것}

\begin{itemize}
\item \textbf{문턱이 표본 셋} --- $0.043$ 과 $0.072$ 사이가 비어 있다.
새 후보가 거기 오면 판정이 바뀔 수 있다.
\item \textbf{덤 --- C 가 모바일을 올린다}($+0.0020$, $11/12$, $0$ 밖
양수). 웹툰의 관측-$0$ 을 결측으로 돌리면 모바일이 조금 오른다 ---
\textbf{이웃 효과}(노트 $551$)다. 판이 내려가므로 안 쫓는다.
\item \textbf{새 자} --- 넷이 다 비었으니 다음 층은 새 검사가 필요하다.
무엇을 볼지는 아직 모른다.
\end{itemize}

\end{document}
